\originalchapter{非线性方程}
\label{ch:nonlinear}

\originalsection{求根}

考虑实变量的非线性关系 $g(x)=h(x)$. 令 $f(x)=g(x)-h(x)$, 就把问题改写为求函数的零点或根, 即 $f(x)=0$. 最基本的假设是函数连续. 下面各方法还会在需要时使用 $f$ 的导数, 并假定我们可以计算函数及所需导数的值.

\textbf{二分法.} 从 $a_0<b_0$ 出发, 假设 $f(a_0)>0$, $f(b_0)<0$. 由连续性, 区间内至少有一个根. 第 $k$ 步已经得到 $[a_k,b_k]$, 且保持相同的异号条件. 令 $m_k=(a_k+b_k)/2$, 将区间分成两半:
\begin{enumerate}
\item 若 $f(m_k)<0$, 保留 $[a_k,m_k]$, 即 $a_{k+1}=a_k$, $b_{k+1}=m_k$.
\item 若 $f(m_k)>0$, 保留 $[m_k,b_k]$, 即 $a_{k+1}=m_k$, $b_{k+1}=b_k$.
\item 若 $f(m_k)=0$, 则 $m_k$ 已是根, 停止算法.
\end{enumerate}
另一种端点符号顺序完全类似. 端点本身为零时也可直接停止.

原文说实际中可在 $f(m_k)$ 足够小时停机. 整理注: 小残差并不单独保证小的位置误差, 尤其在函数很平缓时; 更稳妥的位置容差是看区间宽度. 对嵌套区间选出的极限根 $x^*$, 有
\[
|x^*-m_k|\leq\frac{b_k-a_k}{2}=2^{-(k+1)}(b_0-a_0).
\]
原文给出较宽松的界 $|x^*-m_k|\leq2^{-k}|b_0-a_0|$. 每步区间宽度减半, 相当于增加约一位二进制精度. 二分法的优点是由构造保证收敛到一个根; 缺点是较后面两种方法慢. 若区间内有多个根, 通常不能事先控制算法选择哪一个.

\textbf{Newton--Raphson 法.} 这是一种重要方法, 是科学和工程中许多优化求解器的基础. 先讨论标量方程 $f(x)=0$. 在函数图像的 $(x_n,f(x_n))$ 处作切线, 用切线与横轴的交点作为下一次迭代 $x_{n+1}$. 切线方程给出
\[
0=f(x_n)+(x_{n+1}-x_n)f'(x_n),\qquad
x_{n+1}=x_n-\frac{f(x_n)}{f'(x_n)},
\]
其中要求当前导数不为零.

例如计算 $\sqrt2$, 可求 $f(x)=x^2-2$ 的正根. 迭代为
\[
x_{n+1}=x_n-\frac{x_n^2-2}{2x_n}=\frac{x_n}{2}+\frac1{x_n}.
\]
从 $x_0=1$ 出发, 得 $x_1=3/2=1.5$, $x_2=17/12=1.416666\ldots$, $x_3=577/408=1.4142157\ldots$, 而 $\sqrt2=1.4142135\ldots$.

当 Newton 法收敛时, 速度可以很快. 设 $f''$ 连续, 且根 $x^*$ 附近 $f'$ 不为零; 邻域应包含所考察的迭代点. 记 $E_n=x_n-x^*$. 在 $x_n$ 处展开并在根处求值:
\[
0=f(x^*)=f(x_n)+(x^*-x_n)f'(x_n)+\frac12(x^*-x_n)^2f''(\xi_n),
\]
其中 $\xi_n$ 在 $x_n$ 和 $x^*$ 之间. 与定义下一次迭代的切线式相减, 得
\[
0=-E_{n+1}f'(x_n)+\frac12E_n^2f''(\xi_n),\qquad
E_{n+1}=\frac{f''(\xi_n)}{2f'(x_n)}E_n^2.
\]
若该比值有统一界 $C>0$, 则 $|E_{n+1}|\leq C|E_n|^2$, 迭代后得到
\begin{equation}\label{eq:newton-double-exponential}
|E_n|\leq\frac1C(C|E_0|)^{2^n}.
\end{equation}
称其为二次收敛, 因为下一次误差近似正比于当前误差的平方. 对应的正确数字位数通常约翻倍. 整理注: 原文说“正确数字位数被平方”, 应区分位数翻倍与误差平方. 与二分法的线性收敛相比, 这里的“收敛阶”指迭代误差之间的关系, 与前几章步长误差 $O(h^p)$ 的“方法阶”含义不同.

\begin{remark}
整理补充: 上述误差估计的局部保证不能以尚未证明的 $x_n\to x^*$ 为前提. 可先在闭邻域 $I=[x^*-r,x^*+r]$ 取 $|f'|\geq m>0$, $|f''|\leq M$, 从而取 $C=M/(2m)$; 若 $M=0$, Newton 一步得到根. 取初值满足 $|E_0|\leq r$ 且 $C|E_0|<1$, 则由误差式归纳可知下一点更靠近根, 始终留在 $I$ 中, 式 \eqref{eq:newton-double-exponential} 又趋于零. 这才完成局部收敛保证.
\end{remark}

若初值不够接近, Newton 法可能发散. 导数很小时尤其容易出问题, 因为沿着近水平切线前进可能把点送到很远的区域. 原文给出的例子是 $f(x)=\arctan x$: 初值离原点过远时, Newton 迭代可以发散. 因而它的优势是局部很快, 而一般保证只覆盖根附近.

\textbf{割线法.} 若不知道导数, 可用两个相邻迭代值估计它. 记 $f_n=f(x_n)$, 定义差商
\[
f[x_{n-1},x_n]=\frac{f_n-f_{n-1}}{x_n-x_{n-1}},\qquad
x_{n+1}=x_n-\frac{f_n}{f[x_{n-1},x_n]}.
\]
要求两点不同且分母非零; 已得到根时直接停止. 几何上, 用通过 $(x_{n-1},f_{n-1})$ 与 $(x_n,f_n)$ 的割线代替切线. 它需要两个初始猜测 $x_0,x_1$.

每步只需一次新的函数计算, 因为前一次的 $f_{n-1}$ 已保存. 原文把它与另选很小步长 $h$ 的数值差分近似导数相比: 后者通常要在额外节点计算函数, 而且差分精度还受舍入误差影响.

为了分析收敛, 将割线看成一次插值多项式
\[
p(x)=f_n+f[x_{n-1},x_n](x-x_n).
\]
在两点之间是插值, 在外面是外推. 无论根是否处于两点之间, 只要 $f$ 在包含三个点的区间二阶连续可微, 上一章的余项论证仍给出
\[
f(x)-p(x)=\frac12f''(\xi)(x-x_n)(x-x_{n-1}),
\]
其中 $\xi$ 位于包含 $x,x_{n-1},x_n$ 的最小区间. 在 $x=x^*$ 处代入, 再与 $p(x_{n+1})=0$ 相减, 得
\[
E_{n+1}=\frac{f''(\xi)}{2f[x_{n-1},x_n]}E_nE_{n-1}.
\]
对简单根的足够小邻域, 由中值定理差商等于某点的 $f'$, 因而分母有统一非零下界. 所以
\begin{equation}\label{eq:secant-error}
|E_{n+1}|\leq C|E_n||E_{n-1}|.
\end{equation}
若两个初始误差都足够小, $C\max(|E_0|,|E_1|)<1$, 此式使迭代始终留在邻域并逐步收敛.

原文进一步估算一个介于一阶和二阶之间的收敛阶. 在通常的非退化情形, 假设误差渐近具有 $|E_n|\approx A|E_{n-1}|^p$ 的形式, 与 $|E_{n+1}|\approx K|E_n||E_{n-1}|$ 比较, 左右关于 $|E_{n-1}|$ 的指数分别为 $p^2$ 与 $p+1$. 因而要求 $p^2=p+1$, 正解为
\[
p=\frac{1+\sqrt5}{2}=1.618033\ldots,
\]
即黄金比例. 整理注: 这段是原文的渐近阶推断, 不是从不等式 \eqref{eq:secant-error} 单独推出的等价关系. 当 $f''(x^*)=0$ 等退化情况出现时, 精确的阶可能不同. 通常情形下割线法比二分法快, 比 Newton 的二次收敛慢; 与 Newton 法一样, 一般保证要求初值足够靠近简单根.

\textbf{多元 Newton 法.} 对 $d$ 个未知数和 $d$ 个方程, 设 $f_i(x_1,\ldots,x_d)=0$, $i=1,\ldots,d$, 简记 $f(x)=0$. 方程数与未知数相同只是可能出现孤立解的基本形式, 不保证存在或唯一; 零个、一个或多个解仍须另行分析.

令 $J_f(x)=(\partial f_i/\partial x_j)_{i,j=1}^{d}$ 为 Jacobian 矩阵. 在第 $n$ 次迭代 $x_n$ 附近作向量 Taylor 展开:
\[
0=f(x^*)=f(x_n)+J_f(x_n)(x^*-x_n)+O(\norm{x^*-x_n}^2).
\]
按分量即 $0=f_i(x_n)+\sum_j(\partial f_i/\partial x_j)(x_n)(x_j^*-x_{j,n})+O(\norm{x^*-x_n}^2)$. 整理注: 原文的分量式把这个差的符号颠倒, 此处与向量式统一; $n$ 表示迭代次数, $j$ 表示分量.

忽略二次余项, 定义下一点为
\[
x_{n+1}=x_n-[J_f(x_n)]^{-1}f(x_n).
\]
实际计算应解线性系统 $J_f(x_n)s_n=-f(x_n)$, 再令 $x_{n+1}=x_n+s_n$, 不必显式求矩阵逆. 几何上, 在 $\R^{d+1}$ 中给每个曲面 $y=f_i(x)$ 作切平面, 找到它们交出的直线, 再求与 $y=0$ 的交点. 在 Jacobian 非奇异且足够光滑的简单根附近, 仍有局部二次收敛; 仍要注意初值范围.

\begin{example}\label{ex:circle-sine-newton}
用 Newton 法处理方程组 $x_1^2+x_2^2=1$, $x_2=\sin x_1$.
\end{example}
\begin{solution}
令 $f_1=x_1^2+x_2^2-1$, $f_2=x_2-\sin x_1$. 则
\[
J_f(x)=\begin{pmatrix}2x_1&2x_2\\-\cos x_1&1\end{pmatrix}.
\]
使用二阶矩阵逆公式
\[
\begin{pmatrix}a&b\\c&d\end{pmatrix}^{-1}
=\frac1{ad-bc}\begin{pmatrix}d&-b\\-c&a\end{pmatrix},
\]
得到
\[
J_f(x)^{-1}=\frac1{2x_1+2x_2\cos x_1}
\begin{pmatrix}1&-2x_2\\\cos x_1&2x_1\end{pmatrix},
\]
其中分母须非零. 迭代即
\[
\begin{pmatrix}x_{1,n+1}\\x_{2,n+1}\end{pmatrix}
=\begin{pmatrix}x_{1,n}\\x_{2,n}\end{pmatrix}
-J_f(x_n)^{-1}
\begin{pmatrix}x_{1,n}^2+x_{2,n}^2-1\\x_{2,n}-\sin x_{1,n}\end{pmatrix}.
\]
这给出实际可执行的迭代, 但没有把任意初值的收敛当成已保证.
\end{solution}

\originalsection{优化问题}

科学和工程中的另一类常见问题是寻找函数 $F(x)$ 的最小值或最大值. 若 $x^*$ 的某邻域内所有 $y$ 满足 $F(y)\geq F(x^*)$, 则 $x^*$ 是局部最小点. 若对定义域中全部 $y$ 都如此, 则是全局最小点. 写 $\min_xF(x)$ 表示最小值 $F(x^*)$, 称 $x^*$ 为达到最小值的自变量. 最大化 $F$ 等于最小化 $-F$, 所以下面只谈最小化.

若 $F$ 光滑, 且变量在无约束的实直线上取值, 内点极小值必须满足 $F'(x^*)=0$. 若还检查到 $F''(x^*)>0$, 则得到严格局部最小点的充分条件. 因而可以对 $f=F'$ 应用 Newton 求根法:
\[
x_{n+1}=x_n-\frac{F'(x_n)}{F''(x_n)}.
\]
整理注: $F''(x^*)>0$ 不是所有极小值的必要条件, 如 $x^4$ 在零点极小却二阶导数为零; 驻点也可能是极大点或非极值点.

多元情况下, 最小化标量函数 $F(x_1,\ldots,x_d)$. 无约束内点最小值满足所有一阶偏导为零, 即 $\nabla F(x^*)=0$. 将这看成 $d$ 个非线性方程 $f_i=\partial F/\partial x_i$, 应用 Newton 法得到
\[
x_{n+1}=x_n-[\nabla^2F(x_n)]^{-1}\nabla F(x_n).
\]
矩阵 $\nabla^2F$ 称为 Hessian, 其 $(i,j)$ 元为 $\partial^2F/\partial x_i\partial x_j$. 原文也把它称为 Newton 下降法. 整理注: 只有 Hessian 正定等合适条件下, 完整 Newton 方向才保证是下降方向; 全局求解器通常还需阻尼或线搜索.

与之比较, 简单梯度下降取 $x_{n+1}=x_n-\alpha\nabla F(x_n)$, 其中标量步长 $\alpha$ 足够小. 原文指出梯度下降通常较慢, 但可从更大范围的初值收敛. 这种比较依赖目标函数与步长条件, 不是所有非凸函数上的普遍保证.

\begin{example}\label{ex:log-optimization}
考虑 $F(x_1,x_2)=x_1^2+(\log x_2)^2$, 定义域 $x_1\in\R$, $x_2>0$. 它有唯一的全局最小点.
\end{example}
\begin{solution}
两项均非负, 同时为零恰好在 $(0,1)$, 所以最小值为零. 计算得
\[
\nabla F=\begin{pmatrix}2x_1\\2\log x_2/x_2\end{pmatrix},\qquad
\nabla^2F=\begin{pmatrix}2&0\\0&(2-2\log x_2)/x_2^2\end{pmatrix}.
\]
因此 Newton 迭代为
\[
\begin{pmatrix}x_{1,n+1}\\x_{2,n+1}\end{pmatrix}
=\begin{pmatrix}x_{1,n}\\x_{2,n}\end{pmatrix}
-\begin{pmatrix}1/2&0\\0&x_{2,n}^2/(2-2\log x_{2,n})\end{pmatrix}
\begin{pmatrix}2x_{1,n}\\2\log x_{2,n}/x_{2,n}\end{pmatrix}.
\]
第一分量一步成为零, 因为对二次函数 Newton 法精确地定位极小点. 第二分量可写为 $x_{2,n+1}=x_{2,n}(1-2\log x_{2,n})/(1-\log x_{2,n})$. 整理补充: 在 $x_{2,n}=\ee$ 时 Hessian 奇异, 而某些正初值会被送出 $x_2>0$ 的定义域, 因而这也说明无约束完整 Newton 步的局部性质.
\end{solution}
